% ===================================================================
%  तालिका सं. 2 — मानव शरीर। — अवकाश। — खेल।
%  Traduction 2026 d'apres texto/io, controlee sur texto/fr et
%  texto/en. Consigne : voir texto/hi/10-talika-01.tex.
% ===================================================================

%%K t02-apar-1 apar
\VUsaut{4.0mm}
\VUtitre{15.0pt}{60}{तालिका सं. 2}
\VUsaut{2.0mm}
\VUtitre{11.4pt}{0}{मानव शरीर। --- अवकाश।}
\VUtitre{11.4pt}{0}{खेल।}

%%K t02-tit-0 sub
\VUtitre{13.2pt}{0}{मानव शरीर।}
\VUsaut{2.4mm}
\VUcentre{10.2pt}{0}{\textit{(प्राकृतिक विज्ञान का एक सरल पाठ।)}}
\VUsaut{3.0mm}

%%K t02-01-1 p
1. --- मानव शरीर के मुख्य भाग ये हैं: \VUgras{सिर}, \VUgras{धड़} और
\VUgras{अंग}।

%%K t02-tit-1 sub
\VUpk{10.2pt}{40}{सिर।}
\VUsaut{3.0mm}

%%K t02-02-1 p
2. --- सिर के दो भाग हैं: \VUgras{कपाल}, जिसमें \VUgras{मस्तिष्क} रहता है और
जिसे \VUgras{बाल} ढके रहते हैं, तथा \VUgras{मुख}।

%%K t02-03-1 p
3. --- हमारे चित्र में न कपाल दिखाई देता है न मस्तिष्क; पर मुख के मुख्य भाग
हमें बहुत स्पष्ट दिखाई देते हैं।

%%K t02-04-1 p
4. --- सबसे पहले यह \VUgras{ललाट} है, जो प्रबल बुद्धि और सदा काम करते मन का
पता देता है। \VUgras{कनपटियाँ} बहुत खुली हैं। \VUgras{आँख} चंचल और चौड़ी खुली
है। एक घनी \VUgras{भौंह} और लंबी \VUgras{पलकें} उसकी रक्षा करती हैं।

%%K t02-05-1 p
5. --- यहाँ \VUgras{नाक} और \VUgras{नथुने} भी हैं। नाक से हम सूँघते और साँस
लेते हैं। नाक के दोनों ओर \VUgras{गाल} हैं, और उनसे परे \VUgras{कान}। मुख का
निचला भाग \VUgras{मुँह} और \VUgras{ठोड़ी} से बनता है।

%%K t02-06-1 p
6. --- वे हमें अच्छी तरह दिखाई नहीं देते, क्योंकि \VUgras{मूँछ} और
\VUgras{दाढ़ी} उन्हें छिपाए रखती हैं। पर हम जानते हैं कि मुँह में दो
\VUgras{होंठ}, \VUgras{ऊपरी जबड़ा} और \VUgras{निचला जबड़ा} अपने \VUgras{दाँतों}
समेत, तथा \VUgras{जीभ} होती है। मुँह से हम बोलते और खाते हैं। जीभ से हम अपने
भोजन का स्वाद लेते हैं।

%%K t02-07-1 p
7. --- आँख, कान, नाक और मुँह, उँगलियों के साथ, पाँच इंद्रियों के अंग हैं:
\VUgras{दृष्टि}, \VUgras{श्रवण}, \VUgras{घ्राण}, \VUgras{स्वाद} और
\VUgras{स्पर्श}।

%%K t02-08-1 p
8. --- इस प्रकार सिर मानव शरीर के सबसे रोचक भागों में से एक है।

%%K t02-tit-2 sub
\VUsaut{4.4mm}
\VUpk{10.2pt}{40}{धड़।}
\VUsaut{3.0mm}

%%K t02-09-1 p
9. --- \VUgras{गरदन} सिर को धड़ से जोड़ती है। उसमें \VUgras{गला} और
\VUgras{गुद्दी} आते हैं।

%%K t02-10-1 p
10. --- धड़ में भी बहुत से आवश्यक अंग रहते हैं। \VUgras{छाती} में
\VUgras{फेफड़े} हैं, जिनसे हम साँस लेते हैं, और \VUgras{हृदय}, जो जीवन का
केंद्र है। जब हृदय धड़कना बंद कर देता है, प्राणी मर जाता है।

%%K t02-11-1 p
11. --- \VUgras{पसलियाँ} छाती को सँभालती हैं।

%%K t02-12-1 p
12. --- धड़ के दूसरे भाग हैं \VUgras{पार्श्व}, \VUgras{पीठ}, \VUgras{कंधे} और
\VUgras{उदर} (पेट)। हम सोने के लिए करवट या पीठ के बल लेटते हैं, और बोझ कंधे पर
ढोते हैं।

%%K t02-13-1 p
13. --- उदर में \VUgras{आमाशय} और \VUgras{आँतें} हैं। वे हमें भोजन पचाने में
काम आती हैं।

%%K t02-tit-3 sub
\VUsaut{4.4mm}
\VUpk{10.2pt}{40}{अंग।}
\VUsaut{3.0mm}

%%K t02-14-1 p
14. --- हमारे चार \VUgras{अंग} हैं: दो \VUgras{बाँहें} और दो \VUgras{टाँगें}।
बाँहों से हम वस्तुएँ लेते, पकड़ते, उठाते और बटोरते हैं; टाँगों से हम खड़े होते,
चलते और दौड़ते हैं।

%%K t02-15-1 p
15. --- \VUgras{कंधा} बाँह को शरीर से जोड़ता है। एक बहुत मज़बूत
\VUgras{पेशी}, द्विशिर, बाँह को हिलाती है, जिसे \VUgras{कोहनी}
\VUgras{प्रकोष्ठ} से जोड़ती है। \VUgras{कलाई} \VUgras{हाथ} का जोड़ बनाती है, जो
\VUgras{उँगलियों} और \VUgras{नाखूनों} पर समाप्त होता है। हाथ पर हमें एक
\VUgras{शिरा} (और एक धमनी) दिखाई देती है, जिसमें \VUgras{रक्त} बहता है।

%%K t02-16-1 p
16. --- टाँग के भाग ये हैं: \VUgras{जाँघ}, \VUgras{घुटना} और
\VUgras{घुटने का पिछला भाग}, \VUgras{पिंडली}, जिसका \VUgras{मांस}
\VUgras{त्वचा} से ढका है, \VUgras{टखना} और \VUgras{पैर}। टाँग की
\VUgras{हड्डियाँ} बहुत मोटी और बहुत दृढ़ हैं। पैर के सिरे पर
\VUgras{पैर की उँगलियाँ} हैं। दूसरे भाग \VUgras{तलवा} और \VUgras{एड़ी} हैं।

%%K t02-17-1 p
17. --- मानव शरीर का लगभग पूरा विवरण यही है।

%%K t02-17-2 p
\textit{टिप्पणी।} --- \textit{« मेरा ... दुखता है (मेरा सिर, इत्यादि) » वाक्य
की सहायता से अध्यापक यह सारी शब्दावली अच्छे या बुरे} \VUgras{शारीरिक स्वास्थ्य}
\textit{की दृष्टि से फिर दोहरा सकेंगे।}

%%K t02-tit-4 sub
\VUsaut{5.0mm}
\VUpk{10.2pt}{40}{व्यायाम।}
\VUsaut{2.4mm}
\VUcentre{10.2pt}{0}{\textit{(एक विद्यार्थी की ज़बानी।)}}
\VUsaut{3.0mm}

%%K t02-18-1 p
18. --- लचीले, फुरतीले और स्वस्थ रहने के लिए हमें अपनी टाँगों और बाँहों से काम
लेना चाहिए। इसीलिए हम खेल खेलते हैं और \VUgras{व्यायाम} करते हैं।

%%K t02-19-1 p
19. --- सप्ताह भर ये दोनों अभ्यास विद्यालय के प्रांगण में होते हैं (देखिए
तालिका सं. 1)। पर बृहस्पतिवार और रविवार को हम एक बड़े मैदान में जाते हैं, जहाँ
दौड़ने और मौज करने की जगह मिलती है।

%%K t02-20-1 p
20. --- हमारे अध्यापक और हमारे माता-पिता कभी-कभी साथ आते हैं; पर अभ्यास
\VUgras{व्यायाम-अध्यापक} \textsuperscript{(12)} ही कराते हैं।

%%K t02-21-1 p
21. --- मैदान में, प्रवेश के बाईं ओर, \VUgras{व्यायाम के उपकरण}
\textsuperscript{(1)} खड़े हैं: हम \VUgras{सीढ़ी} \textsuperscript{(2)} पर चढ़ते हैं,
\VUgras{छल्लों} \textsuperscript{(3)} और \VUgras{झूला-डंडे} \textsuperscript{(4)} पर
उलटे झूलते हैं, \VUgras{रस्सी की सीढ़ी} \textsuperscript{(5)} या
\VUgras{गाँठदार रस्सी} \textsuperscript{(6)} के सिरे तक हाथों-हाथ चढ़ जाते हैं।

%%K t02-22-1 p
22. --- इस बीच हमारे साथी \VUgras{काठ के घोड़े} \textsuperscript{(7)} या
\VUgras{समांतर डंडों} \textsuperscript{(11)} के ऊपर से कूदते हैं। कुछ
\VUgras{मुक्केबाज़ी} \textsuperscript{(8)} करते हैं या मुखौटा और \VUgras{फ़ॉयल}
लेकर \VUgras{तलवारबाज़ी} \textsuperscript{(9)} करते हैं। कुछ अंत में
\VUgras{डंबल} \textsuperscript{(10)} उठाते हैं।

%%K t02-tit-5 sub
\VUsaut{4.6mm}
\VUpk{10.2pt}{40}{खेल।}
\VUsaut{3.0mm}

%%K t02-23-1 p
23. --- व्यायाम करने वालों के चारों ओर लड़के और लड़कियाँ तरह-तरह के
\VUgras{खेल} खेल रहे हैं। लड़कियाँ अपने \VUgras{घेरे} \textsuperscript{(14)} से मन
बहलाती हैं; वे \VUgras{रस्सी कूदती} \textsuperscript{(13)} हैं, या अपने भाइयों और
चचेरे भाइयों को \VUgras{क्रोके} \textsuperscript{(15)} के खेल में बुलाती हैं। जब
विरोधी सोचता है कि वह सहज ही घेरे में से निकल जाएगा, तभी वे अपने
\VUgras{काठ के मुँगरे} से उसकी \VUgras{गेंद} दूर मार देती हैं, और इससे उन्हें
बड़ा आनंद आता है!

%%K t02-24-1 p
24. --- लड़के अधिक ज़ोर के खेल पसंद करते हैं। वे प्रसन्नता से
\VUgras{नौगोलियाँ} \textsuperscript{(16)} खेलते हैं, जिन्हें वे \VUgras{गेंद}
\textsuperscript{(17)} से साफ़ गिरा देते हैं। वे \VUgras{लट्टू} \textsuperscript{(18)}
और \VUgras{बिलबोके} \textsuperscript{(19)} को, यहाँ तक कि
\VUgras{मेंढक के खेल} \textsuperscript{(20)} को भी, तुच्छ नहीं समझते। पर उनके
प्रिय खेल \VUgras{कंचे} \textsuperscript{(21)} और \VUgras{दीवार-गेंद}
\textsuperscript{(22)} हैं, क्योंकि \VUgras{काँच-घर} \textsuperscript{(26)} की दीवार
हलकी गेंद उछालने के लिए ठीक बैठती है।

%%K t02-25-1 p
25. --- नन्हा जॉन अपने \VUgras{बाँसों} \textsuperscript{(24)} पर चढ़ा बहुत कुशल है
और अपना संतुलन खूब सँभालता है। उसके पास कुछ साथी उसी काँच-घर में और
\VUgras{बाड़} के पीछे \VUgras{लुका-छिपी} \textsuperscript{(25)} खेल रहे हैं। कुछ
को \VUgras{हथेली का खेल} \textsuperscript{(28)} अधिक भाता है, या
\VUgras{चिड़िया} \textsuperscript{(29)}, जो एक \VUgras{बल्ले} से दूसरे तक उड़ती
रहती है।

%%K t02-noto-1 noto
\VUnotes{143.2mm}{%
(*) बच्चों के खेलों के नाम प्रायः बिलकुल राष्ट्रीय होते हैं। हमने उन्हें इदो
में जैसे बन पड़ा वैसे नाम दिए हैं — कभी उस क्रिया से जो उन्हें पहचान देती है,
कभी किसी न किसी देश का नाम अपनाकर, जब वह बहुत भ्रामक न हो। उदाहरण के लिए:
\VUgras{पीपे का खेल} (जर्मन और फ़्रेंच) वही \VUgras{मेंढक का खेल}
\textsuperscript{(20)} (अंग्रेज़ी) है, जिसमें खिलाड़ी अपनी गोल चकतियाँ छिद्रों
वाले संदूक (पीपे) पर मुँह खोले खड़े धातु के मेंढक के मुँह में फेंकने का यत्न
करते हैं। \VUgras{कंधा-कूद} \textsuperscript{(30)} \VUgras{बकरी-कूद} से केवल इतना
भिन्न है: सिर के ऊपर से कूदने के लिए खिलाड़ी अपने साथी के कंधों पर हाथ रखते
हैं, जब कि « \VUgras{बकरी-कूद} » में वे केवल उसकी पीठ पर हाथ रखते हैं। --- या
फिर, कुछ खिलाड़ी झुक जाते हैं और दूसरे उनकी पीठ पर हाथ रखकर उनके ऊपर से कूदते
हैं; पहलों को बिना डिगे धक्का सहना होता है, दूसरों को बिना संतुलन खोए कूदना
(तालिका ही देखिए)।%
}

%%K t02-26-1 p
26. --- मैदान के बीच दो पेड़ हैं। उनमें से एक के नीचे सात-आठ लड़कों ने
\VUgras{कंधा-कूद} \textsuperscript{(30)} (*) का खेल जमाया है। यह खेल हर देश में
नहीं जाना जाता; पर \VUgras{बकरी-कूद} क्या है, यह सब जानते हैं, जिसे बच्चे इस
समय नहीं खेल रहे।

%%K t02-tit-6 sub
\VUsaut{5.0mm}
\VUpk{10.2pt}{40}{खुली हवा के खेल।}
\VUsaut{3.4mm}

%%K t02-27-1 p
27. --- \VUgras{आँख-मिचौली} \textsuperscript{(31)} में भी बड़ा मज़ा आता है;
लड़कियाँ और लड़के बारी-बारी से आँखों पर \VUgras{पट्टी} बाँधते हैं और उन
फूहड़ों को पकड़ते हैं जो पकड़े जाने देते हैं।

%%K t02-28-1 p
28. --- मैदान के बीच, यह रहा एक बहुत फ़्रेंच खेल, यद्यपि कुछ खुरदरा: यह
\VUgras{रीछ} \textsuperscript{(32)} है, जो \VUgras{पतंग} \textsuperscript{(33)} के
शांत खेल से, यहाँ तक कि अंग्रेज़ी खेल \VUgras{टेनिस} \textsuperscript{(34)} से भी,
कहीं अधिक शोरगुल वाला है।

%%K t02-29-1 p
29. --- वह अंतिम खेल बड़े विद्यार्थियों का आनंद है। हमारे अध्यापक स्वयं भी उसे
प्रसन्नता से अपनाते हैं। उसमें बड़ी कुशलता और फुरती चाहिए, और वह मुझे बहुत
भाता है। \VUgras{झूला} \textsuperscript{(35)} मैं लड़कियों के लिए छोड़ देता हूँ।

%%K t02-30-1 p
30. --- एक और अंग्रेज़ी खेल भी मुझे नहीं भाता, जो सहज ही ख़तरनाक हो सकता है।

%%K t02-30-1b p
मेरा अभिप्राय \VUgras{फ़ुटबॉल} \textsuperscript{(36)} से है, जो मेरे बड़े भाई का
आनंद है। मुझे हमारा पुराना खेल \VUgras{बंदी-अड्डा} \textsuperscript{(38)} अधिक
भाता है, उतना ही स्वास्थ्यकर और कहीं कम कठोर।

%%K t02-tit-7 sub
\VUsaut{4.6mm}
\VUpk{10.2pt}{40}{साइकिल और गेंद के खेल।}
\VUsaut{3.0mm}

%%K t02-31-1 p
31. --- मुझे मैदान के चारों ओर \VUgras{साइकिल} \textsuperscript{(37)} चलाना भी
अच्छा लगता है।

%%K t02-32-1 p
32. --- अगले वर्ष मैं \VUgras{घुड़सवारी} \textsuperscript{(39)} सीखूँगा। घोड़ा
कैसा सुंदर लगता है जब वह ठुमकता या सुडौल दुलकी चलता है, और सरपट बाधाएँ पार कर
जाता है!

%%K t02-33-1 p
33. --- गरमियों में हम \VUgras{तैरना} \textsuperscript{(40)} सीखते हैं। हम नदी के
स्वच्छ ठंडे पानी में आनंद से डुबकी लगाते और नहाते हैं। कभी-कभी अंत में हम एक
कागज़ का \VUgras{गुब्बारा} \textsuperscript{(41)} उड़ाते हैं, जो गरम हवा से भरते
ही गंभीरता से ऊपर उठ जाता है।

%%K t02-34-1 p
34. --- और भी बहुत से खेल हैं, पर उन्हें गिनाते-गिनाते मैं कभी न चुकूँगा, और
इस छोटे से वृत्तांत से ही दिखाई देता है कि हमारे सुंदर मैदान पर बृहस्पतिवार
नीरस होने से बहुत दूर हैं।

%%K t02-34-2 p
न. ब. --- \textit{अध्यापक इस अध्याय को सहज ही पूरा कर सकेंगे, प्रत्येक
महत्त्वपूर्ण खेल का अधिक विस्तार से वर्णन करके।}
\VUsaut{6.0mm}
\VUfilet{22mm}
